% Generated draft: requires mathematical and source review.
\documentclass[11pt]{article}
\usepackage[margin=25mm]{geometry}
\usepackage{fontspec}
\setmainfont{Latin Modern Roman}
\usepackage{amsmath,amssymb,mathtools,unicode-math}
\setmathfont{Latin Modern Math}
\usepackage{enumitem,needspace,xurl,hyperref}
\setlength{\parindent}{0pt}
\setlength{\parskip}{5pt}
\setlength{\emergencystretch}{3em}
\newcommand{\N}{\mathbb{N}}
\newcommand{\Z}{\mathbb{Z}}
\newcommand{\Q}{\mathbb{Q}}
\newcommand{\R}{\mathbb{R}}
\newcommand{\C}{\mathbb{C}}
\newcommand{\F}{\mathbb{F}}
\newcommand{\A}{\mathbb{A}}
\newcommand{\PP}{\mathbb{P}}
\newcommand{\E}{\mathbb{E}}
\newcommand{\eps}{\varepsilon}
\newcommand{\dd}{\,\mathrm{d}}
\begin{document}
\section*{UnsolvedMath ID: 30005444}
\textbf{Source problem code:} catalogue entry
\subsection{Definitions and mathematical statement}
The supplied record asks about a planar extension of the Bricmont--Kupiainen invariance principle. Let \(\mathcal C_{\mathrm{BK}}(2)\) denote the class of random-environment models satisfying the intended Bricmont--Kupiainen hypotheses in dimension \(d=2\), and let \(\mathsf{IP}_{\mathrm{full}}(M)\) denote the full invariance principle intended in the source for a model \(M\). These symbols are labels by reference: the supplied record does not state the hypotheses defining this class or the precise convergence assertion defining this property.

The source question has the logical form
\[
\text{Does }\quad
\forall M\in\mathcal C_{\mathrm{BK}}(2),\qquad
\mathsf{IP}_{\mathrm{full}}(M)
\quad\text{hold?}
\]
This is a formulation template requiring completion from the cited primary sources. In particular, the record does not specify the transition mechanism, environment law, independence or mixing assumptions, ellipticity, perturbation size, symmetry assumptions, starting conditions, centering, scaling, path-space topology, or limiting covariance. The supplied background describes the higher-dimensional result as a perturbative quenched invariance principle, but does not provide a complete definition of the requested planar conclusion. A self-contained mathematical statement cannot therefore be recovered without further source evidence.

The original extract also mentions a weaker symmetry condition than that of Bricmont--Kupiainen and work in progress of Bolthausen and an unidentified author that would prove an assertion labeled \((1)\) for \(d=2\). Neither that assertion nor the weaker symmetry condition is supplied. Their relationship to the displayed question is unresolved; they cannot be treated as established results or incorporated as definite hypotheses.
\subsection{Short English statement}
The question asks whether the full invariance principle under the Bricmont--Kupiainen assumptions extends to dimension \(d=2\). The supplied record omits the assumptions and the exact convergence statement, so this draft preserves the target by reference and explicitly records the missing information. The designation as open is supplied dated curation, not an independently verified conclusion.
\subsection{Sources}
\begin{itemize}
\item Local UnsolvedMath registry, permanent ID \(30005444\), title \emph{Planar Bricmont--Kupiainen Invariance Principle}, source problem code \texttt{catalogue entry}. The record belongs to the Oberwolfach Reports Open Problems collection and gives the proposed year as 2023; no proposer is identified.
\item \emph{Mini-Workshop: New Horizons in Motions in Random Media} (2023), Oberwolfach Reports, pp. 565--606, DOI 10.4171/owr/2023/10. Supplied original-report reference: \url{https://doi.org/10.4171/owr/2023/10}. The original extract includes a page marker for p. 573 and incomplete surrounding text.
\item Bricmont--Kupiainen, \emph{Random walks in asymmetric random environments}. Supplied literature reference: \url{https://doi.org/10.1007/BF02102067}.
\item Bolthausen--Zeitouni, described in the supplied record as \emph{multiscale exit distributions}. Supplied literature reference: \url{https://arxiv.org/abs/math/0607192}.
\item Supplied literature assessment, checked on \texttt{2026-08-22}: the record labels the problem open, attributes a perturbative quenched invariance principle for \(d\geq3\) to Bricmont--Kupiainen, and states that the report leaves the full \(d=2\) case as a problem and that no primary source located in that curation closes it. This is supplied dated curation, not independently verified here; its time-sensitive status and cited claims require independent checking.
\item The registry labels the statement \texttt{corrected\_verified} and reports that the official Oberwolfach report was checked directly and that the clean statement preserves the complete mathematical target while restoring definitions and removing extraction artifacts or source-dependent wording. These are supplied verification claims; this draft has not independently visited or checked the cited sources, and the supplied clean statement does not actually include the missing definitions.
\end{itemize}
\end{document}
